\section{Introduction}

A neural code records which members of a finite family of receptive fields
are active at each point of an ambient space.  When the fields are required
to be open and convex, the resulting combinatorial objects link convex
geometry, combinatorics, topology, and mathematical neuroscience
\citep{CurtoItskovVelizYoungs2013,CurtoEtAl2017,CruzEtAl2019}.  Even in the
plane, a realization may begin with curved, tangent, or degenerate
boundaries, so a basic structural question is whether the same code can be
represented by polygons.

\citet{BukhJeffs2023} answered this affirmatively: every planar open convex
code has a realization by interiors of convex polygons.  Their proof first
passes to a suitable closed realization, then chooses an inclusion-minimal
subrealization and proves that every member of such a minimal tuple is
polygonal.  The theorem is exact at the level of the code, but its statement
does not require a specified cross-section or a specified atom witness to
remain fixed.

The purpose of this paper is to show that the planar construction has an
exact relative form.  Finitely many line sections can be frozen before the
minimal shrinking step, and finitely many selected atom witnesses can be
protected.  The key finite datum on a prescribed line is not merely the
order of endpoints after clipping by the ambient polygon.  One records the
complete bounded line section first.  A four-point trace frame forces every
point of its open section to remain interior after shrinking, while the
closed endpoints prevent movement of the section.  Only afterward is the
line section clipped by the ambient polygon.  This order of operations is
what preserves half-open traces at ambient vertices and simultaneous endpoint
events.

Our main result is stated in \cref{thm:main}.  In addition to the open trace
identity requested there, the proof preserves the complete closed section
of every prescribed line after an auxiliary bounded reduction.  This
slightly stronger formulation makes tangent contacts and point sections
transparent.  Protected witnesses are retained by small triangles placed
inside their active fields.

The relative theorem has two immediate interface applications.  First, if
two planar face problems share an edge and inherit the same source trace,
their polygonal outputs agree pointwise in strict membership on that edge.
Second, all six edge traces of a source tetrahedron can be prescribed
simultaneously on its four triangular faces, producing a coherent finite
polygonal model of the complete boundary sphere.  These statements concern
membership traces only.  They do not match numerical defining functions,
normal slopes, or three-dimensional convex fillings.

\paragraph{Relationship to prior work.}
The baseline polygonalization theorem and the local simplification mechanism
are due to \citet{BukhJeffs2023}.  The new ingredient is a finite trace frame
and the resulting relative bookkeeping: arbitrary finitely many line
sections, protected atom witnesses, a fixed ambient polygon, tangent closed
contacts, compatible shared-edge outputs, and the tetrahedral-boundary
application.  Standard approximation of convex bodies by polytopes concerns
metric error rather than exact Boolean atoms and exact prescribed sections;
see, for example, \citet{SchneiderWieacker1981}.  Exact interpolation also
appears in the opposite smoothing direction---for example, prescribed
subsets of polytope facets can be retained while smoothing a single body
\citep{Ghomi2004}---but that problem does not preserve a multi-set neural
atom code.  Our argument is instead an exact constrained shrinking argument
for an entire convex family.

\paragraph{Organization.}
After definitions in \cref{sec:prelim}, \cref{sec:background} reviews the
inclusion-minimal planar construction in the form needed here.
\Cref{sec:traces} builds the finite trace data and proves the main theorem.
\Cref{sec:degenerate} treats tangent, clipped, and lower-dimensional cases.
\Cref{sec:interfaces} gives shared-edge and finite facewise compatibility,
and \cref{sec:tetra} proves the tetrahedral-boundary application.
\Cref{sec:limitations} records the precise scope and the absence of a
three-dimensional filling theorem.
